% Part: computability
% Chapter: recursive-functions
% Section: other-recursions

\documentclass[../../../include/open-logic-section]{subfiles}

\begin{document}

\olfileid[pt]{cmp}{rec}{ore}
\olsection{Outras formas de recursão}

Usando pareamento e sequenciamento, podemos justificar formas mais
exóticas (e úteis) de recursão primitiva. Por exemplo, muitas vezes é
útil definir duas funções simultaneamente, como na seguinte definição:
\begin{align*}
h_0(\vec x, 0) & = f_0(\vec x) \\
h_1(\vec x, 0) & = f_1(\vec x) \\
h_0(\vec x, y+1) & = g_0(\vec x, y, h_0(\vec x, y), h_1(\vec x, y)) \\
h_1(\vec x, y+1) & = g_1(\vec x, y, h_0(\vec x, y), h_1(\vec x, y))
\end{align*}
Esse é um caso de \emph{recursão simultânea}. Outra maneira útil de
definir funções é dar o valor de $h(\vec x, y+1)$ em termos de
\emph{todos} os valores $h(\vec x, 0)$, \dots,~$h(\vec x, y)$, como
na seguinte definição:
\begin{align*}
h(\vec x, 0) & = f(\vec x) \\
h(\vec x, y+1) & = g(\vec x, y, \tuple{h(\vec x, 0), \dots, h(\vec x, y)}).
\end{align*}
O esquema a seguir expressa essa ideia de forma mais sucinta:
\[
h(\vec x, y) = g(\vec x, y, \tuple{h(\vec x, 0), \dots, h(\vec x, y-1)})
\]
com o entendimento de que o último argumento de $g$ é simplesmente a
sequência vazia quando $y$ é $0$. Em qualquer uma das formulações, a
ideia é que, ao calcular o ``passo sucessor'', a função $h$ pode usar
toda a sequência de valores calculados até então.
Isso é conhecido como \emph{recursão por curso de valores}. Como exemplo
específico, ela pode ser usada para justificar o seguinte tipo de
definição:
\begin{align*}
h(\vec x, y) & = \begin{cases}
  g(\vec x, y, h(\vec x, k(\vec x, y))) & \text{se $k(\vec x, y) < y$} \\
  f(\vec x) & \text{caso contrário}
\end{cases}
\end{align*}
Em outras palavras, o valor de $h$ em $y$ pode ser calculado em termos
do valor de $h$ em \emph{qualquer} valor anterior, dado por~$k$.


\begin{prob}
  Defina a função resto $r(x,y)$ por recursão por curso de valores.
  (Se $x$, $y$ são números naturais e $y > 0$, $r(x,y)$ é o número
  menor que~$y$ tal que $x = z\times y + r(x,y)$ para algum~$z$.
  Para fixar a definição, digamos que, se $y=0$, $r(x,0) = 0$.)
\end{prob}

Você deve pensar em como obter essas funções usando recursão primitiva
comum. Uma última versão da recursão primitiva é mais flexível, pois
permite alterar os \emph{parâmetros} (valores laterais) ao longo do
processo:
\begin{align*}
h(\vec x, 0) & = f(\vec x) \\
h(\vec x, y+1) & = g(\vec x, y, h(k(\vec x), y))
\end{align*}
Isso também pode ser simulado pela recursão primitiva comum. (Fazer
isso é trabalhoso. Como dica, tente desenrolar o cálculo à mão.)

\end{document}

